% ==========================================================================
% GENERATED by scripts/split_paper.py. Do not edit this file.
%
% Source: imported_papers/probe-norm/content/body.tex, content/supplement.tex
% Part 4: Results
%
% Edits belong upstream, in the published paper. To change what the
% chapter includes, change imports/bounded.toml and re-run
% the splitter.
% ==========================================================================
\section{Results}
\label{sec:exp}
\subsection{Synthetic latents}
\label{sec:res_suite}
\begin{figure}[t]
\centering
\includegraphics[width=\linewidth]{figures/fig_suite.pdf}
\caption{\textbf{Probes fail on synthetic latents whose norm drifts, and re-projection repairs them.} \textbf{(A)} Error of the ordinary probe, ReLU and linear, as the norm of a latent that stores the state in its direction grows or shrinks while the state stays the same; the re-projected probe returns to its error at the norm of the encodings. A norm that grows by $5$\% per prediction step reaches ten times its value after $47$ steps. \textbf{(B)} Share of pairs of latents that the ordinary and the re-projected probe order by their true state error, when the norm varies independently of that error; points are seeds, marks their means.}
\label{fig:suite}
\end{figure}
We first build latents whose content we control, to test whether re-projection resolves the problems of the ordinary probe identified in Section~\ref{sec:split}. Each latent stores a physical state in its direction, so a probe that reads the state correctly gives the same answer at any norm. We then change the norm and leave the state as it is, which is what autoregressive prediction does when it compounds a small error in norm. An error of $5$\% per step multiplies the norm by ten within $47$ steps. At that norm the ordinary probe's error is $27$ times that of always predicting the mean state, although the latent still encodes the state exactly, and the re-projected probe returns to its error at the original norm (Figure~\ref{fig:suite}A). When we give each latent a known state error and a norm chosen independently of it, the ordinary probe orders only about $61$\% of pairs of latents by their true error, where chance is 50\%, and the re-projected probe orders all of them (Figure~\ref{fig:suite}B). If the state is stored in the norm instead, re-projection would remove it, and the radius check detects this case before any correction (\supsec~\ref{app:synthetic}). On synthetic predictions for which we know which error is correct, or that neither is, the protocol decides correctly in $50$ of $50$ cases (\supsec~\ref{app:protocol_test}). Re-projection therefore removes the effect of the norm on both the error and the ranking of these latents, and the radius and support checks flag the cases in which it would mislead.

\subsection{Trained world models}
\label{sec:res_drift}
\begin{figure}[t]
\centering
\includegraphics[width=\linewidth]{figures/fig_horizon.pdf}
\caption{Share of the probe's error that the norm causes in the trained modules under \textbf{(A)} the cosine and \textbf{(B)} the squared-error loss. Left axis: re-projection factor, the error after re-projection over the error before. Right axis: the share of the error that re-projection removes. Medians over runs and tasks. Shaded: trained horizon ($T\le H=6$).}
\label{fig:horizon}
\end{figure}

\input{tables/table_arms.tex}

\paragraph{Free norms.}
Under cosine loss, twenty steps ahead, the ordinary probe reports an error $185$ times larger than that of the re-projected probe for the modules that leave the norm free (Table~\ref{tab:arms}). Per run, re-projection removes a median $99$\% of the ordinary probe's error (Figure~\ref{fig:horizon}A). Without the encoder's final normalization, the ordinary probe's error is $179$ times the re-projected one. With a shared encoder and actions that move the state, it is $1.3\times10^{4}$ times larger. Linear and ReLU probes respond to the norm in this way, while a tanh probe saturates, so its error stays bounded however far the norm drifts (\suptab~\ref{tab:v2_capacity}).

\paragraph{Rankings.}
The two probes also rank the modules differently, so they disagree on which model is better (Table~\ref{tab:arms}). Under cosine loss, Orthogonal comes first under the ordinary probe and last under the re-projected probe, because it keeps its norm while the other modules drift. Twenty steps ahead, the two probes disagree on four of the six module pairs under cosine loss, on five without final normalization, and on all six with a shared encoder. Sections~\ref{sec:res_rank} and~\ref{sec:res_checks} test whether such rankings hold, and on models of known accuracy re-projection ranks correctly where the ordinary probe fails.

\paragraph{Controlled norms.}
Where the design keeps the norm in check, both probes agree. Rescaling every predicted step to the initial norm gives $2.13$ under the ordinary probe and $2.11$ after re-projection, SimNorm $1.39$ and $1.46$, and in both settings the two probes rank the modules alike, up to one swap of neighbors. Squared-error training keeps norms close within the training horizon. There re-projection raises the error for ten of the twelve module--task combinations, because shorter predictions have less accurate directions, as expected when squared error averages over uncertain futures. The norm then carries the model's uncertainty, and the protocol keeps the raw error. Beyond the training horizon drift returns, with $9.2$ against $6.3$ twenty steps ahead (Figure~\ref{fig:horizon}B).

\subsection{Published world models}
\label{sec:res_lewm}
\input{tables/table_published.tex}

For TCRL on pendulum (Table~\ref{tab:published}), the ordinary probe reports $0.11$ of the mean state's error at the end of the training horizon, and re-projection lowers it to $2.3\times10^{-3}$. Assigning the predicted norms to the true target encodings reproduces the ordinary probe's error, so the predicted directions are accurate and the norm causes the ordinary probe's error (Figure~\ref{fig:intro}B; \supsec~\ref{app:problem}). For TCRL on cheetah-run, re-projection lowers the error from $9.1\times10^{3}$ to $1.2$ times the mean state's, but random directions with the same norms improve just as much, so the protocol attributes the reduction to the probe's sensitivity to scale. Random directions carry no state, so an improvement that they share says nothing about the predictions. For the Koopman autoencoders and for LeWorldModel on Two-Room, re-projection also lowers the error, but most predictions lie away from the encodings, and the protocol calls the reading inconclusive. A probe or decoder fitted on encodings extrapolates on such predictions, so neither its raw nor its re-projected error can be trusted there. For TD-MPC, TD-MPC2, LeWorldModel on Push-T and Semigroup-JEPA, both probes agree and no correction is needed. Each of these evaluations controls the norm. TD-MPC2 sets it with SimNorm, LeWorldModel bounds it with a normalized output head, TD-MPC is evaluated at the end of its training horizon, where squared error keeps it close, and the Semigroup-JEPA probe normalizes its input. KSL, PlayVirtual and LightZero apply their loss to the output of a projection network instead of the latent itself, so training never compares their raw predicted latents with the encodings. In raw latent space none of their predictions is supported, and the protocol calls the reading inconclusive. In the projection space, where the loss acts, the re-projected probe beats persistence and the mean state five steps ahead on every checkpoint (\suptab~\ref{tab:published_more}).

\subsection{Rankings with a known answer}
\label{sec:res_rank}
\input{tables/table_known.tex}

To test rankings against a known truth, we build models of known accuracy from the encoders of our trained models (Section~\ref{sec:protocol}). Each such model predicts the encoding of a real frame whose state lies close to the true future state, and the closer that frame, the more accurate the model. In every pair, the more accurate model also has the larger norm drift. The ordinary probe ranks the worse model first in all $120$ of these pairs (Table~\ref{tab:known}). For the $30$ pairs of equally accurate models, it declares one of the two better every time. The protocol confirms the true ranking in $88$ pairs, including all $30$ at the largest accuracy gap. For the other $32$, where the two models differ least, it reports the ranking as inconclusive, and it never confirms a wrong one. These pairs are most frequent on cartpole and acrobot, where the probe's error is larger even on the true encodings, so small differences between models fall within it.

\subsection{What the checks add}
\label{sec:res_checks}
\label{sec:consequences}
A lower re-projected error means that the norm distorted the score, but the model may still predict poorly. The checks detect two such cases in our trained modules. On pendulum and acrobot, re-projection makes the rollouts look accurate, but the start-encoding probe does at least as well (\suptab~\ref{tab:controls}). A ranking of these modules then says nothing about their dynamics, and the protocol rejects it. On cartpole the rollouts beat every null control, but the ordinary probe is inaccurate on their directions, while a probe fitted on the predictions has a much lower error. The support check flags the predictions that lie away from the encodings, and the protocol calls the cartpole rankings inconclusive. The same holds with a shared encoder up to fifty steps ahead, and no choice of thresholds confirms any of these rankings (\suptabs~\ref{tab:verdicts} and~\ref{tab:sensitivity}).
%%END_RESULTS
